% First the desired computation, then the authenticated route around the square.
\begin{frame}{Background: Evaluation over QAS}
  \centering
  \begin{tikzpicture}[
      >={Latex[length=2.8mm]},
      state/.style={inner sep=5pt,font=\large,text=CVQCInk},
      process/.style={draw=CVQCInk,fill=CVQCInk,rounded corners=3mm,
        minimum width=1.7cm,minimum height=0.85cm,
        text=white,font=\large},
      flow/.style={->,very thick,draw=CVQCTeal},
      label/.style={font=\large,text=CVQCInk}
    ]
    \path[use as bounding box] (-6.6,-3.15) rectangle (6.6,3.05);

    \node[state] (input) at (-3.6,1.8) {$\ket{\psi}$};
    \node[state] (output) at (3.6,1.8) {$C\ket{\psi}$};
    \draw[flow] (input.east) -- node[label,above=5pt] {$C$} (output.west);

    \only<2>{
      \node[process] (enc) at (-3.6,0) {$\operatorname{Enc}$};
      \node[process] (dec) at (3.6,0) {$\operatorname{Dec}$};
      \node[state] (cipherinput) at (-3.6,-1.8)
        {$\operatorname{Enc}(\ket{\psi})$};
      \node[state] (cipheroutput) at (3.6,-1.8)
        {$\operatorname{Enc}(C\ket{\psi})$};
      \node[state] (reject) at (5.8,0) {$\mathsf{rej}$};

      \draw[flow] (input.south) -- (enc.north);
      \draw[flow] (enc.south) -- (cipherinput.north);
      \draw[flow] (cipherinput.east) --
        node[label,below=6pt] {$\operatorname{Eval}(C)$} (cipheroutput.west);
      \draw[flow] (cipheroutput.north) -- (dec.south);
      \draw[flow] (dec.north) -- (output.south);
      \draw[flow] (dec.east) -- (reject.west);
    }
  \end{tikzpicture}

  \note<1>{
    \footnotesize
    We want to apply a circuit C to the authenticated input. The top arrow
    is the desired logical computation. The pure-state, unitary notation is
    a cartoon: the authenticated computation also has to preserve correlations
    with an external reference, and our eventual application has measurements.
  }
  \note<2>{
    \footnotesize
    The lower route performs the computation while the data stays
    authenticated. For honest prescribed evaluation, decoding gives the same
    logical result as the top arrow. Enc on the right uses the updated key;
    we suppress keys throughout the picture. The verifier tracks those updates
    without revealing the key to the prover. Dec includes verification and
    can reject; the branch on the right recalls the security goal against
    deviations: the intended result or rejection, with negligible total
    accepted error. This is not an arbitrary QAS feature: we need a scheme
    and gate procedures that support authenticated evaluation. Some gates
    need interaction and auxiliary authenticated states. It is also distinct
    from the classical-client QFHE used later in the remote state preparation
    step. The next slide illustrates the particularly simple CNOT procedure
    for the trap code. The manuscript's Non-CCZ Evaluation section describes
    the operation and the verifier's key updates.
  }
\end{frame}
